% GENERATED FILE. Edit canonical lessons, then run npm run generate:books.
% canonical-source-hash: fnv1a64:2407b0146f180aeb
% canonical-lessons: IT-C296-sospettato, IT-C296-movente, IT-C296-alibi, IT-C296-interrogare, IT-C296-confessare, IT-R296-ripasso-289-296

\chapter{Chi è stato}
\label{ch:it-chi-e-stato}
\hlchaptermodality{296}

\begin{chapteropening}
\textbf{By the end of this chapter:} \emph{I can follow a whodunnit: the suspect, the motive, the alibi, the questioning and the confession.}

Complete the last lesson of chapter 296: I can follow a whodunnit: the suspect, the motive, the alibi, the questioning and the confession.
\end{chapteropening}

\section[il sospettato]{il sospettato — the suspect}

\label{lesson:IT-C296-sospettato}



\begin{quote}
Before the new one: say the Italian for the clue, then the Italian for the footprint.
\end{quote}

\subsection*{You'll want to know: il sospettato}
\textbf{il sospettato} — \textquotedblleft{}the suspect\textquotedblright{}.

\begin{quote}
Il sospettato ha quarant'anni.
\end{quote}

\begin{quote}
\emph{The suspect is forty.}
\end{quote}

From \emph{sospettare}, ‘to suspect’.

\begin{cousinweb}
\emph{sospettato} is the past participle of \emph{sospettare} \textquotedblleft{}to suspect\textquotedblright{}, built on \emph{sospetto}, from Latin \emph{suspectus} \textquotedblleft{}looked at from below\textquotedblright{}, from \emph{suspicere} (Lewis \& Short).
\end{cousinweb}

\begin{grammarlens}[title={What you've built}]
One more word for telling, explaining and describing at B1.
\end{grammarlens}

\subsection*{Your turn}
\begin{itemize}
\raggedright
  \item \emph{Say it:} \emph{il sospettato}
  \item \emph{Say it:} \emph{il sospettato}, once more
  \item \emph{Say it:} a sentence of your own with \emph{il sospettato}
  \item \emph{Recall:} say the Italian for the clue, then the Italian for the footprint, then say \emph{il sospettato} again
\end{itemize}

\begin{tcolorbox}[breakable,colback=teal!4,colframe=teal!35!black,title={Before you move on}]
What does il sospettato mean? (\textquotedblleft{}The suspect\textquotedblright{}.) Say it once more.
\end{tcolorbox}
\section[il movente]{il movente — the motive (for a crime)}

\label{lesson:IT-C296-movente}



\begin{quote}
Before the new one: say the Italian for the footprint, then the Italian for the suspect.
\end{quote}

\subsection*{You'll want to know: il movente}
\textbf{il movente} — \textquotedblleft{}the motive (for a crime)\textquotedblright{}.

\begin{quote}
Il movente del delitto è il denaro.
\end{quote}

\begin{quote}
\emph{The motive for the crime is money.}
\end{quote}

For crimes; in daily life a reason is \emph{il motivo}.

\begin{cousinweb}
\emph{movente} is the present participle of \emph{muovere} \textquotedblleft{}to move\textquotedblright{}, from Latin \emph{movēre} (Lewis \& Short): what moves someone to act.
\end{cousinweb}

\begin{grammarlens}[title={What you've built}]
One more word for telling, explaining and describing at B1.
\end{grammarlens}

\subsection*{Your turn}
\begin{itemize}
\raggedright
  \item \emph{Say it:} \emph{il movente}
  \item \emph{Say it:} \emph{il movente}, once more
  \item \emph{Say it:} a sentence of your own with \emph{il movente}
  \item \emph{Recall:} say the Italian for the footprint, then the Italian for the suspect, then say \emph{il movente} again
\end{itemize}

\begin{tcolorbox}[breakable,colback=teal!4,colframe=teal!35!black,title={Before you move on}]
What does il movente mean? (\textquotedblleft{}The motive (for a crime)\textquotedblright{}.) Say it once more.
\end{tcolorbox}
\section[l'alibi]{l'alibi — the alibi}

\label{lesson:IT-C296-alibi}



\begin{quote}
Before the new one: say the Italian for the suspect, then the Italian for the motive.
\end{quote}

\subsection*{You'll want to know: l'alibi}
\textbf{l'alibi} — \textquotedblleft{}the alibi\textquotedblright{}.

\begin{quote}
Lui ha un alibi: era al cinema.
\end{quote}

\begin{quote}
\emph{He has an alibi: he was at the cinema.}
\end{quote}

Plural unchanged: \emph{gli alibi}.

\begin{cousinweb}
\emph{alibi} was taken from Latin \emph{alibī} \textquotedblleft{}elsewhere\textquotedblright{} (EtymDB), built on \emph{alius} \textquotedblleft{}other\textquotedblright{} and \emph{ibī} \textquotedblleft{}there\textquotedblright{}.
\end{cousinweb}

\begin{grammarlens}[title={What you've built}]
One more word for telling, explaining and describing at B1.
\end{grammarlens}

\subsection*{Your turn}
\begin{itemize}
\raggedright
  \item \emph{Say it:} \emph{l'alibi}
  \item \emph{Say it:} \emph{l'alibi}, once more
  \item \emph{Say it:} a sentence of your own with \emph{l'alibi}
  \item \emph{Recall:} say the Italian for the suspect, then the Italian for the motive, then say \emph{l'alibi} again
\end{itemize}

\begin{tcolorbox}[breakable,colback=teal!4,colframe=teal!35!black,title={Before you move on}]
What does l'alibi mean? (\textquotedblleft{}The alibi\textquotedblright{}.) Say it once more.
\end{tcolorbox}
\section[interrogare]{interrogare — to question; to test orally (at school)}

\label{lesson:IT-C296-interrogare}



\begin{quote}
Before the new one: say the Italian for the motive, then the Italian for the alibi.
\end{quote}

\subsection*{You'll want to know: interrogare}
\textbf{interrogare} — \textquotedblleft{}to question; to test orally (at school)\textquotedblright{}.

\begin{quote}
La polizia ha interrogato il sospettato per ore.
\end{quote}

\begin{quote}
\emph{The police questioned the suspect for hours.}
\end{quote}

At school, \emph{interrogare} is to test a pupil out loud.

\begin{cousinweb}
\emph{interrogare} is a learned borrowing from Latin \emph{interrogāre} \textquotedblleft{}to ask, to question\textquotedblright{} (Lewis \& Short), from \emph{inter} and \emph{rogāre} \textquotedblleft{}to ask\textquotedblright{}.
\end{cousinweb}

\begin{grammarlens}[title={What you've built}]
One more word for telling, explaining and describing at B1.
\end{grammarlens}

\subsection*{Your turn}
\begin{itemize}
\raggedright
  \item \emph{Say it:} \emph{interrogare}
  \item \emph{Say it:} \emph{interrogare}, once more
  \item \emph{Say it:} a sentence of your own with \emph{interrogare}
  \item \emph{Recall:} say the Italian for the motive, then the Italian for the alibi, then say \emph{interrogare} again
\end{itemize}

\begin{tcolorbox}[breakable,colback=teal!4,colframe=teal!35!black,title={Before you move on}]
What does interrogare mean? (\textquotedblleft{}To question; to test orally (at school)\textquotedblright{}.) Say it once more.
\end{tcolorbox}
\section[confessare]{confessare — to confess}

\label{lesson:IT-C296-confessare}



\begin{quote}
Before the new one: say the Italian for the alibi, then the Italian for to question.
\end{quote}

\subsection*{You'll want to know: confessare}
\textbf{confessare} — \textquotedblleft{}to confess\textquotedblright{}.

\begin{quote}
Alla fine ha confessato tutto.
\end{quote}

\begin{quote}
\emph{In the end he confessed everything.}
\end{quote}

\emph{Confessarsi} is to go to confession in church.

\begin{cousinweb}
\emph{confessare} is built on Latin \emph{cōnfessus}, the past participle of \emph{cōnfitērī} \textquotedblleft{}to admit\textquotedblright{} (Lewis \& Short).
\end{cousinweb}

\begin{grammarlens}[title={What you've built}]
One more word for telling, explaining and describing at B1.
\end{grammarlens}

\subsection*{Your turn}
\begin{itemize}
\raggedright
  \item \emph{Say it:} \emph{confessare}
  \item \emph{Say it:} \emph{confessare}, once more
  \item \emph{Say it:} a sentence of your own with \emph{confessare}
  \item \emph{Recall:} say the Italian for the alibi, then the Italian for to question, then say \emph{confessare} again
\end{itemize}

\begin{tcolorbox}[breakable,colback=teal!4,colframe=teal!35!black,title={Before you move on}]
What does confessare mean? (\textquotedblleft{}To confess\textquotedblright{}.) Say it once more.
\end{tcolorbox}
\section[(dialogue)]{Ripasso dei capitoli 289–296}

\label{lesson:IT-R296-ripasso-289-296}



\begin{quote}
Say the Italian for the reader and the listener.
\end{quote}

\subsection*{Your turn}
Say these aloud:

\begin{itemize}
\raggedright
  \item the reader, the listener, the viewer, the audience, the editorial
  \item the crime, the serious crime, the criminal, the victim, to commit
  \item the bag-snatching, to snatch, the pickpocket, to rob someone, the stolen goods
  \item the robbery, to rob, the robber, the safe, the loot
  \item the break-in, to force open, the burglar alarm, the alarm, the vandal
  \item the carabiniere, the police station, the patrol, the video camera, the chase
  \item to investigate, the police investigation, the inspector, the clue, the footprint
  \item the suspect, the motive, the alibi, to question, to confess
\end{itemize}

\begin{tcolorbox}[breakable,colback=teal!4,colframe=teal!35!black,title={Before you move on}]
The reader? (\textbf{il lettore}.) The listener? (\textbf{l'ascoltatore}.) The viewer? (\textbf{lo spettatore}.) The audience? (\textbf{il pubblico}.) The editorial? (\textbf{l'editoriale}.) The crime? (\textbf{il crimine}.) The serious crime? (\textbf{il delitto}.) The criminal? (\textbf{il criminale}.) The victim? (\textbf{la vittima}.) To commit? (\textbf{commettere}.) The bag-snatching? (\textbf{lo scippo}.) To snatch? (\textbf{scippare}.) The pickpocket? (\textbf{il borseggiatore}.) To rob someone? (\textbf{derubare}.) The stolen goods? (\textbf{la refurtiva}.) The robbery? (\textbf{la rapina}.) To rob? (\textbf{rapinare}.) The robber? (\textbf{il rapinatore}.) The safe? (\textbf{la cassaforte}.) The loot? (\textbf{il bottino}.) The break-in? (\textbf{lo scasso}.) To force open? (\textbf{scassinare}.) The burglar alarm? (\textbf{l'antifurto}.) The alarm? (\textbf{l'allarme}.) The vandal? (\textbf{il vandalo}.) The carabiniere? (\textbf{il carabiniere}.) The police station? (\textbf{il commissariato}.) The patrol? (\textbf{la pattuglia}.) The video camera? (\textbf{la telecamera}.) The chase? (\textbf{l'inseguimento}.) To investigate? (\textbf{indagare}.) The police investigation? (\textbf{l'indagine}.) The inspector? (\textbf{l'ispettore}.) The clue? (\textbf{l'indizio}.) The footprint? (\textbf{l'impronta}.) The suspect? (\textbf{il sospettato}.) The motive? (\textbf{il movente}.) The alibi? (\textbf{l'alibi}.) To question? (\textbf{interrogare}.) To confess? (\textbf{confessare}.)
\end{tcolorbox}
